%&pdflatex
%%%%%%%%%%%%%%%%%%%%%%%%%%%%%%%%%%%%%%%%%%%%%%%%%%%%%%%%%%%%%%%%%%%%%%
%                                                                    %
% Mean and Weak Convergence of Some Orthogonal Fourier Expansions    % 
% by Using $A_{p}$ Theory,                                           %
% by                                                                 %
% J. J. Guadalupe, M. P\'erez and J. L. Varona                       %
%                                                                    %
% PUBLISHED IN:                                                      %
% ``Orthogonal Polynomials and Their Applications''                  %
% (Proc. Int. Congr., Laredo 1987, J. Vinuesa, ed.),                 %
% 161-169, Lect. Notes Pure Appl. Math. 117,                         %
% Marcel Dekker, New York, 1989.                                     %
% ISBN: 0-8247-8161-9                                                %
%                                                                    %
% Rewrited with LaTeX2e: 2-September-2004                            %
%                                                                    %
%%%%%%%%%%%%%%%%%%%%%%%%%%%%%%%%%%%%%%%%%%%%%%%%%%%%%%%%%%%%%%%%%%%%%%

\documentclass[11pt,a4paper]{article}

\nofiles % The file .aux is not necessary

\usepackage{amsmath,amsfonts,amsthm}

\usepackage{indentfirst}

\addtolength{\textwidth}{0.5cm}

\newcommand{\R}{\mathbb{R}}
\newcommand{\N}{\mathbb{N}}

\newcommand{\supp}{\operatorname{supp}}

\makeatletter
\renewcommand\section{\@startsection {section}{1}{\z@}%
                                   {-3.5ex \@plus -1ex \@minus -.2ex}%
                                   {2.3ex \@plus.2ex}%
                                   {\centering \normalfont\large\bfseries}}
\makeatother

\theoremstyle{plain}
\newtheorem*{theoMNT}{Theorem}
\newtheorem*{lemOne}{Lema 1}
\newtheorem*{lemTwo}{Lema 2}
\newtheorem*{lemThree}{Lema 3}
\newtheorem*{lemFour}{Lema 4}
\newtheorem*{theoOne}{Theorem 1}
\newtheorem*{theoTwo}{Theorem 2}
\newtheorem*{theoThree}{Theorem 3}
\newtheorem*{corOne}{Corollary 1}

\theoremstyle{definition}
\newtheorem*{definition}{Definition}

\theoremstyle{remark}
\newtheorem*{remOne}{Remark 1}
\newtheorem*{remTwo}{Remark 2}
\newtheorem*{remThree}{Remark 3}

\title{Mean and Weak Convergence\\ of Some Orthogonal Fourier Expansions\\ 
by Using $A_{p}$ Theory}

\author{\scshape Jos\'e J. Guadalupe, 
Mario P\'erez, Juan L. Varona\\[5pt]
\textit{Departamento de Matem\'aticas,}
\textit{Universidad de Zaragoza, Spain}%
\renewcommand\thefootnote{}%
\renewcommand\footnotemark{}%
\footnote{Authors are supported in part 
by a grant of the CAICYT, PB85-0338.}%
\footnote{THIS PAPER HAS BEEN PUBLISHED IN:
\textit{Lect.\ Notes Pure Appl.\ Math.\@}
\textbf{117} (1989), 161--169.}%
\footnote{Copyright \copyright\ 1989 by Marcel Dekker, Inc.}%
}

\date{}

%%%%%%%%%%%%%%%%%%%%%%%%%%%%%%%%%%%%%%%%%%%%%%%%%%%%%%
\begin{document}
%%%%%%%%%%%%%%%%%%%%%%%%%%%%%%%%%%%%%%%%%%%%%%%%%%%%%%

%%%%%%%%%%%%%%%%%%%%%%%%%%%%%%%%%%%%%%%%%%%%%%%%%%%%%%
\maketitle
%%%%%%%%%%%%%%%%%%%%%%%%%%%%%%%%%%%%%%%%%%%%%%%%%%%%%%


%%%%%%%%%%%%%%%%%%%%%%%%%%%%%%%%%%%%%%%%%%%%%%%%%%%%%%
\section*{Introduction}
%%%%%%%%%%%%%%%%%%%%%%%%%%%%%%%%%%%%%%%%%%%%%%%%%%%%%%

Let $d\mu$ be a finite positive Borel measure on $\R$ such that 
$\supp(d\mu)$ is an infinite set and let $p_{n}(d\mu)$ denote 
the corresponding orthonormal polynomials. For $f \in L^{1}(d\mu)$, 
$S_{n}f$ stands for the $n$th partial sum of the orthogonal Fourier 
expansion of $f$ in $\{p_{n}(d\mu)\}_{n=0}^{\infty}$, 
that is,
$$
S_{n}(f,x) = \sum_{k=0}^{n} a_{k}p_{k}(x),
\quad a_{k} = \widehat{f}(k) 
= \int_{\R} fp_{k}\,d\mu.
$$

The study of the convergence of $S_{n}f$ in $L^{p}(d\mu)$ ($p \neq 2$)
has been discussed for several classes of orthogonal polynomials 
(c.f. Askey-Wainger [1], Badkov [2--4], Muckenhoupt [9--11], Newman-Rudin 
[13], Pollard [14--16], Wing [19]). For instance, in the case 
of Jacobi polynomials $\{P^{(\alpha,\beta)}_{n}(x)\}_{n=0}^{\infty}$ 
which are orthogonal in $[-1,1]$ with respect to the weight $w(x) 
= (1-x)^{\alpha}(1+x)^{\beta}$, $\alpha, \beta \geq -1/2$,
Pollard proved that
$|1/p-1/2| < \min\{1/(4\alpha+4),1/(4\beta+4)\}$ 
is a sufficient condition for the uniform boundedness
$\|S_{n}f\|_{p,w} \leq C\|f\|_{p,w}$, 
which is equivalent to the convergence in $L^{p}(w)$, $1 < p < \infty$. 
Newman and Rudin showed that the previous condition is also necessary 
and later Muckenhoupt extended these results to $\alpha, \beta > -1$.

On the other hand, M\'at\'e, Nevai and Totik [8] obtained, in 
a general way, necessary conditions for the mean convergence 
of Fourier expansions:

%% \textsc{Theorem (Mate-Nevai-Totik)}.
\begin{theoMNT}[M\'at\'e-Nevai-Totik]
Let $d\mu$ be such that $\supp(d\mu) = [-1,1]$,
$\mu' > 0$ almost everywhere, $U$ and $V$ nonnegative 
Borel measurable functions such that neither of them vanishes 
almost everywhere in $[-1,1]$ and $V$ is finite on a set with positive 
Lebesgue mea-sure. If $S_{n}$ is uniformly bounded from $L^{p}(V^{p}\,d\mu)$ 
into $L^{p}(U^{p}\,d\mu)$, then
\begin{itemize}
\item[(i)] $U^{p} \in L^{1}(d\mu)$, $V^{-q} \in L^{1}(d\mu)$, 
$q = p/(p-1)$,
\item[(ii)] $\displaystyle \int_{-1}^{1} U(x)^{p} \mu'(x)^{1-p/2}
(1-x^{2})^{-p/4}\,dx < \infty$,
\item[(iii)] $\displaystyle \int_{-1}^{1} V(x)^{-q} \mu'(x)^{1-q/2}
(1-x^{2})^{-q/4}\,dx < \infty$.
\end{itemize}
\end{theoMNT}

%%%%%%%%%%%%%%%%%%%%%%%%%%%%%%%%%%%%%%%%%%%%%%%%%%%%%%
\section*{Mean convergence}
%%%%%%%%%%%%%%%%%%%%%%%%%%%%%%%%%%%%%%%%%%%%%%%%%%%%%%

The main subject in this paper is the study of the mean and weak 
boundedness of the orthogonal Fourier expansion, in some particular 
cases, by using $A_{p}$-theory, which plays a central role in the 
weighted norm inequalities for the Hardy-Littlewood maximal operator 
and the Hilbert transform.

We start with the mean convergence recalling some definitions:
\begin{itemize}
\item[(i)] $(u,v) \in A_{p}(-1,1)$, $1 < p < \infty$, iff there 
exists a positive constant $C$ such that
$\displaystyle \left( \int_{I} u(x)\,dx \right)
\left( \int_{I} v(x)^{-1/(p-1)}\,dx \right)^{p-1} \leq C|I|^{p}$
for all intervals $I \subset [-1,1]$, where $|I|$ is the 
Lebesgue measure of~$I$.

\item[(ii)] $(u,v) \in A_{p}^{\delta}(-1,1)$ ($\delta > 1$)
iff $(u^{\delta},v^{\delta}) \in A_{p}(-1,1)$.

\item[(iii)] \sloppy Given a sequence $\{(u_{n},v_{n})\}_{n\in\N}$, we say that 
$(u_{n},v_{n}) \in A_{p}(-1,1)$ uniformly if there exists a constant 
$C$, independent of $n$, such that
$\displaystyle \left( \int_{I} u_{n}(x)\,dx \right)
\left( \int_{I} v_{n}(x)^{-1/(p-1)}\,dx \right)^{p-1} \leq C|I|^{p}$ 
for all intervals $I \subset [-1,1]$.
\end{itemize}

It is well known that $(u,v) \in A_{p}$ is a necessary condition 
for the boundedness of the Hilbert transform $H$ from $L^{p}(v)$ into 
$L^{p}(u)$ and that $(u,v) \in A_{p}^{\delta}$ (for some 
$\delta > 1$) is a sufficient condition [7], [12]. Analogous 
conditions work for the uniform boundedness, modifying slightly 
the arguments in~[12].

This is connected with the Fourier expansion, and the idea comes 
from Pollard: let $\{p_{n}(x)\}_{n=0}^{\infty}$ denote 
the orthonormal polynomials with respect to $d\mu = \mu'(x)\,dx$ 
and $\{q_{n}(x)\}_{n=0}^{\infty}$ the orthonormal 
polynomials with respect to $(1-x^{2})\,d\mu$. Then
$$
S_{n}(f,x) = \int_{-1}^{1} f(t)K_{n}(x,t)\mu'(t)\,dt
$$
and the kernel $K_{n}(x,t)$ can be decomposed in the form
$$
K_{n}(x,t) = r_{n}T_{1}(n,x,t) + s_{n}T_{2}(n,x,t) + 
s_{n}T_{3}(n,x,t)
$$
where:
\begin{align*}
T_{1}(n,x,t) &= p_{n+1}(x)p_{n+1}(t),\\
T_{2}(n,x,t) &= (1-t^{2}) \frac{p_{n+1}(x) q_{n}(t)}{x-t},\\
T_{3}(n,x,t) &= T_{2}(n,t,x) = (1-x^{2}) \frac{p_{n+1}(t) q_{n}(x)}{t-x}.
\end{align*}
If $\mu' > 0$ a.e., then $\{r_{n}\}$ and $\{s_{n}\}$ are bounded~[17].
Let $U$ and $V$ be weights, $1 < p < \infty$, and
$$
W_{i}(f,x) = W_{i,n}(f,x) = \int_{-1}^{1} f(t)T_{i}(n,x,t)\mu'(t)\,dt
\quad (i = 1,2,3).
$$
We try to estimate the three terms:
$$
\|(W_{i}f)U\|_{p,\mu'} \leq C\|fV\|_{p,\mu'}.
$$
Denote:
\begin{gather*}
u_{n}(x) = |p_{n+1}(x)|^{p}U(x)^{p}\mu'(x), \quad
v_{n}(x) = |q_{n}(x)|^{-p}(1-x^{2})^{-p}V(x)^{p}\mu'(x)^{1-p},\\
\overline{u}_{n}(x) = |q_{n}(x)|^{p}(1-x^{2})^{p}U(x)^{p}\mu'(x), \quad
\overline{v}_{n}(x) = |p_{n+1}(x)|^{-p}V(x)^{p}\mu'(x)^{1-p}.
\end{gather*}
By using H\"{o}lder's inequality and $A_{p}$ results we obtain the 
following sufficient conditions for the boundedness of $W_{i}$ ($i = 1,2,3$):
\begin{align*}
(u_{n},v_{n}) &\in A_{p}^{\delta}(-1,1) \quad\text{uniformly for 
some } \delta > 1,\\
(\overline{u}_{n},\overline{v}_{n}) &\in A_{p}^{\delta}(-1,1) 
\quad\text{uniformly for some } \delta > 1.
\end{align*}
On the other hand, the conditions
\begin{equation}
((1-x^{2})^{-p/4}U(x)^{p}\mu'(x)^{1-p/2},
(1-x^{2})^{-p/4}V(x)^{p}\mu'(x)^{1-p/2}) \in A_{p}(-1,1)
\tag{1}
\end{equation}
and
\begin{equation}
((1-x^{2})^{p/4}U(x)^{p}\mu'(x)^{1-p/2},
(1-x^{2})^{p/4}V(x)^{p}\mu'(x)^{1-p/2}) 
\in A_{p}(-1,1)
\tag{2}
\end{equation}
turn out to be necessary for the boundedness of $W_{i}$ ($i = 1,2,3$). 
From (1), (2) and Th.~2 in~[8], M\'at\'e-Nevai-Totik's conditions 
for the mean convergence of $S_{n}f$ can be obtained.

Next, we introduce a particular kind of measures.

\begin{definition}
We say that $d\mu = \mu'(x)\,dx \in H$ 
if $\mu'(x) = (1-x)^{\alpha}(1+x)^{\beta}w(x)$, where:
\begin{itemize}
\item[(i)] $w > 0$ a.e.\ and $C_{1} < w(x) < C_{2}$ for $x \in (1-\varepsilon,1)$ 
and $x \in (-1,-1+\varepsilon)$.
\item[(ii)] $|p_{n}(x)| \leq C (1 - x + a_{n})^{-(\alpha/2+1/4)}
(1 + x + b_{n})^{-(\beta/2+1/4)}w(x)^{-1/2}$.
\item[(iii)] $|q_{n}(x)| \leq C (1 - x + a_{n})^{-(\alpha/2+3/4)}
(1 + x + b_{n})^{-(\beta/2+3/4)} w(x)^{-1/2}$
and $\{a_{n}\}$, $\{b_{n}\}$ are positive sequences such that $\lim a_{n} 
= \lim b_{n} = 0$.
\end{itemize}
\end{definition}

There exist particular weights belonging to the class~$H$: the 
generalized Jacobi weights $(GJ)$ $d\mu(x) = \mu'(x)\,dx$, 
being
$$
\mu'(x) = \varphi(x)(1-x)^{\Gamma_{1}}
\prod_{k=2}^{N-1} |x-x_{k}|^{\Gamma_{k}} (1+x)^{\Gamma_{N}}
$$
where $\Gamma_{k} \geq 0$ ($k = 1,2,\dots,N$),
$1 > x_{2} > \dots  > x_{N-1} > -1$, 
$\varphi > 0$ and continuous on $[-1,1]$ and
$\omega(\delta)/\delta \in L^{1}(0,1)$, $\omega$
being the modulus of continuity of~$\varphi$.

\begin{theoOne}
Let $d\mu \in H$, 
$U(x)=(1-x)^{a}(1+x)^{b}u(x)$, 
$V(x)=(1-x)^{A}(1+x)^{B}v(x)$, with 
$u > 0$ a.e., $v > 0$ a.e.\ and such that $C_{1} < u(x), v(x) < C_{2}$ 
for $x \in (1-\varepsilon,1)$ and $x \in (-1,-1+\varepsilon)$. 
If
\begin{align*}
|(\alpha+1)(1/p-1/2) + (a+A)/2| 
&< (a-A)/2 + \min\{1/4,(\alpha+1)/2\},\quad A \leq a,\\
|(\beta+1)(1/p-1/2) + (b+B)/2| 
&< (b-B)/2 + \min\{1/4,(\beta+1)/2\},\quad B \leq b,
\end{align*}
and
$$
  (w^{1-p/2}u^{p},w^{1-p/2}v^{p}) \in A_{p}^{\delta}(-1,1) 
  \quad\text{for some \ } \delta > 1,
$$
then:
$$
\int_{-1}^{1} |S_{n}(f,x)U(x)|^{p}\mu'(x)\,dx 
\leq C\int_{-1}^{1} |f(x)V(x)|^{p}\mu'(x)\,dx.
$$
\end{theoOne}

This theorem is a consequence of the following lemmas:

\begin{lemOne}
Let $\{u_{n}(x)\}$, $\{v_{n}(x)\}$, $\{U_{n}(x)\}$, $\{V_{n}(x)\}$ 
be sequences of weights defined on a finite interval $(a,b)$. Let 
$c \in (a,b)$ and $\varepsilon > 0$ be fixed and independent 
of~$n$. Assume that there exist some positive constants $\lambda_{i}$ 
($i = 1,2,3,4$) such that $\lambda_{1} \leq U_{n}(x), V_{n}(x) 
\leq \lambda_{2}$ on $(a,c+\varepsilon)$ and $\lambda_{3} 
\leq u_{n}(x), v_{n}(x) \leq \lambda_{4}$ on $(c-\varepsilon,b)$.\\
If $(u_{n},v_{n}) \in A_{p}(a,c)$ and $(U_{n},V_{n}) \in A_{p}(c,b)$ 
uniformly, then $(u_{n}U_{n},v_{n}V_{n}) \in A_{p}(a,b)$ 
uniformly.
\end{lemOne}

\begin{lemTwo}
Let $\{x_{n}\}$ be a sequence of positive numbers which 
converges to zero. Then
$(x^{r}(x+x_{n})^{s},x^{R}(x+x_{n})^{S}) \in A_{p}^{\delta}(0,1)$ 
uniformly if and only if:
$$
r > -1,\ R < p-1,\ R \leq r,\ R+S \leq r+s,\ r+s > -1,\ R+S < p-1.
$$
\end{lemTwo}

From the above theorem, we have the following result, which was 
established by Badkov [3] (using other methods and without the 
restriction $\Gamma_{k} \geq 0$, $2 \leq k \leq N-1$):

\begin{corOne}
Let $w \in (GJ)$ and
$U(x) = (1-x)^{a}(1+x)^{b} \prod_{k=2}^{N-1} |x-x_{k}|^{c_{k}}$. 
If
\begin{align*}
|(\Gamma_{1}+1) (1/2-1/p) - a| < \min\{1/4,(\Gamma_{1}+1)/2\},\\
|(\Gamma_{N}+1) (1/2-1/p) - b| < \min\{1/4,(\Gamma_{N}+1)/2\}
\end{align*}
and
$$
|(\Gamma_{k}+1) (1/2-1/p) - c_{k}| < \min\{1/2,(\Gamma_{k}+1)/2\}
\quad (k = 2,\dots ,N-1),
$$
then
$$
\int_{-1}^{1} |S_{n}(f,x)U(x)|^{p}\mu'(x)\,dx 
\leq C \int_{-1}^{1} |f(x)U(x)|^{p}\mu'(x)\,dx.
$$
\end{corOne}


%%%%%%%%%%%%%%%%%%%%%%%%%%%%%%%%%%%%%%%%%%%%%%%%%%%%%%
\section*{Weak convergence}
%%%%%%%%%%%%%%%%%%%%%%%%%%%%%%%%%%%%%%%%%%%%%%%%%%%%%%

Another aim in this paper is to examine the weak behaviour of 
the orthogonal Fourier expansion, that is to study if there exists 
a constant $C$, independent of $n$, $y$ and $f$, such that:
$$
\int_{|S_{n}(f,x)|>y} d\mu(x)
\leq Cy^{-p} \int_{-1}^{1} |f(x)|^{p}\,d\mu(x), 
\quad y > 0,
$$
i.e., if $S_{n}$ is uniformly bounded from $L^{p}(d\mu)$ into 
$L^{p}_{*}(d\mu)$, $1 < p < \infty$.

The previous inequality only can be true, besides the mean convergence 
interval, in its endpoints. For the Fourier-Legendre expansion 
($d\mu = dx$), Chanillo [5] proved that the partial sum operator 
is not weak type $(4,4)$.

The following result gives necessary conditions for the weak 
boundedness~[6].

\begin{theoTwo}
Let $d\mu$ be such that $\supp(d\mu) = [-1,1]$, 
$\mu' > 0$ a.e., $U$ and $V$ be weights, $1 < p < \infty$. 
If there exists a constant $C$ such that
$$
\|S_{n}f\|_{L^{p}_{*}(U^{p}\,d\mu)} \leq C\|f\|_{L^{p}(V^{p}\,d\mu)}
$$
holds for all integers $n \geq 0$ and every $f \in L^{p}(V^{p}\,d\mu)$, 
then:
\begin{itemize}
\item[(i)] $U^{p}, V^{-q} \in L^{1}(d\mu)$,
\item[(ii)] $\mu'(x)^{-1/2}(1-x^{2})^{-1/4} \in L^{p}_{*}(U^{p}\mu'\,dx)$,
\item[(iii)] $\mu'(x)^{-1/2}(1-x^{2})^{-1/4} \in L^{q}(V^{-q}\mu'\,dx)$.
\end{itemize}
\end{theoTwo}

This result is a consequence of the following lemmas:

\begin{lemThree}
Let $U$ and $V$ be weights and $1 < p < \infty$. 
If there exists a constant $C$ such that for every $f \in L^{p}(V^{p}\,d\mu)$ 
the inequality
$$
\|S_{n}f\|_{L^{p}_{*}(U^{p}\,d\mu)} \leq C\|f\|_{L^{p}(V^{p}\,d\mu)}
$$
holds for all integers $n \geq 0$, then
$$
\|p_{n}\|_{L^{q}(V^{-q}\,d\mu)} \|p_{n}\|_{L^{p}_{*}(U^{p}\,d\mu)} \leq C.
$$
\end{lemThree}

\begin{lemFour}[{[8, Th.~2]}]
Let $\supp(d\mu) = [-1,1]$, $\mu' > 0$ a.e.\ in
$[-1,1]$ and $0 < p \leq \infty$. There exists a
constant $C$ such that if $g$ is a Lebesgue-measurable function 
in $[-1,1]$, then
$$
\|\mu'(x)^{-1/2}(1-x^{2})^{-1/4}\|_{L^{p}(|g|^{p}\,dx)}
\leq \liminf_{n \to \infty} \|p_{n}\|_{L^{p}(|g|^{p}\,dx)}.
$$
In particular, if
$$
\liminf_{n \to \infty} \|p_{n}\|_{L^{p}(|g|^{p}\,dx)} = 0
$$
then $g = 0$ a.e.
\end{lemFour}

We are going to study the weak boundedness of the Fourier-Jacobi 
expansion. Since for $-1 < \alpha, \beta \leq -1/2$ 
the conditions $|1/p-1/2|< \min\{1/(4\alpha+4),1/(4\beta+4)\}$ 
are trivial for $p \in (1,\infty)$, we suppose, by symmetry, 
$\alpha \geq \beta$ and $\alpha > -1/2$. Then, 
the mean convergence interval is $4(\alpha+1)/(2\alpha+3) 
< p < 4(\alpha+1)/(2\alpha+1)$.

\begin{remOne}
If $U(x) = V(x) = 1$, the inequality $(\alpha+1)(1/p-1/2) 
< 1/4$ is not satisfied for $p = 4(\alpha+1)/(2\alpha+3)$. 
It implies that $S_{n}$ is not weak type $(p,p)$ for the lower endpoint 
of the interval of mean convergence. The same happens with generalized 
Jacobi polynomials.
\end{remOne}

\begin{remTwo}
The conditions in Theorem~2 are the same as those 
of M\'at\'e-Nevai-Totik's theorem. Thus, the conditions obtained 
by M\'at\'e, Nevai and Totik are necessary not only for the mean 
convergence but also for the weak convergence.
\end{remTwo}

\begin{remThree}
It can be proved that M\'at\'e-Nevai-Totik's conditions 
are not sufficient for the weak convergence. In order to prove 
this, consider the Fourier-Legendre expansion ($d\mu = dx$), 
$p = 4$ and take
\begin{align*}
U(x) &= \left|\log \frac{1+x}{4} \right|^{-5/8}
\left|\log \frac{1-x}{4} \right|^{-5/8}, \\
V(x) &= \left|\log \frac{1+x}{4} \right|^{-3/8}
\left|\log \frac{1-x}{4} \right|^{-3/8}.
\end{align*}
\end{remThree}

Let $S_{n}$ denote the $n$th partial sum of the Fourier-Jacobi expansion 
with respect to $\mu'(x) = (1-x)^{\alpha}(1+x)^{\beta}$, 
being $\alpha \geq \beta$ and $\alpha > -1/2$. 
Then, the interval of mean convergence is given by $4(\alpha+1)/(2\alpha+3) 
< p < 4(\alpha+1)/(2\alpha+1)$. Theorem~2 works 
to prove that $S_{n}$ is not weak type on $L^{p}(\mu')$ for $p 
= 4(\alpha+1)/(2\alpha+3)$, but it is not useful to show 
that $S_{n}$ is not weak type for $p = 4(\alpha+1)/(2\alpha+1)$. 
It leads us to make use of other arguments.

\begin{theoThree}
Let $r = 4(\alpha+1)/(2\alpha+1)$. Then, there exists no 
constant $C$, independent of $n$ and $f \in L^{r}(\mu')$, such that
$$
\|S_{n}f\|_{L^{r}_{*}(\mu')} \leq C\|f\|_{L^{r}(\mu')}.
$$
\end{theoThree}

\begin{proof}
Decompose the kernel $K_{n}(x,t)$, as before, in the 
form
$$
K_{n}(x,t) = r_{n}T_{1}(n,x,t) + s_{n}T_{2}(n,x,t) + 
s_{n}T_{3}(n,x,t).
$$
By using the estimates
$$
|p_{n}(x)| \leq C(1-x)^{-\alpha/2-1/4}, 
\quad
|q_{n}(x)| \leq C(1-x)^{-\alpha/2-3/4}, 
\quad x \in (0,1),
$$
H\"{o}lder's inequality and standard arguments of $A_{p}$ theory, 
the boundedness of $T_{1}$ and $T_{3}$ can be proved.

Now, it is not difficult to prove that
$$
\int_{|p_{n}(x)H(f(t)q_{n-1}(t)(1-t^{2})\mu'(t),x)|>y} \mu'(x)\,dx
\leq Cy^{-r} \|f\|^{r}_{r,\mu'}
$$
is not satisfied for any fixed constant $C$. The proof is by contradiction, 
constructing a sequence of functions $\{f_{m,n}\}$ such that the 
constant $C$ appearing in the previous inequality grows with~$m$.
\end{proof}

%%%%%%%%%%%%%%%%%%%%%%%%%%%%%%%%%%%%%%%%%%%%%%%%%%%%%%
\begin{thebibliography}{99}
%%%%%%%%%%%%%%%%%%%%%%%%%%%%%%%%%%%%%%%%%%%%%%%%%%%%%%

\frenchspacing

\bibitem{1.}
\textsc{R. Askey and S. Wainger}, Mean convergence of expansions 
in Laguerre and Hermite series, \textit{Amer. J. Math.} \textbf{87} (1965), 
695--708.

\bibitem{2.}
\textsc{V. M. Badkov}, Convergence in the mean of Fourier series 
in orthogonal polynomials, \textit{Math. Notes} \textbf{14} (1973), 651--657.

\bibitem{3.}
\textsc{V. M. Badkov}, Convergence in the mean and almost everywhere 
of Fourier series in polynomials orthogonal on an interval, 
\textit{Math. USSR Sbornik} \textbf{24} (1974), 223--256.

\bibitem{4.}
\textsc{V. M. Badkov}, Approximation properties of Fourier series 
in orthogonal polynomials, \textit{Russian Math. Surveys} \textbf{33} (1978), 
53--117.

\bibitem{5.}
\textsc{S. Chanillo}, On the weak behaviour of partial 
sums of Legendre series, \textit{Trans. Amer. Math. Soc.} \textbf{268} 
(1981), 367--376.

\bibitem{6.}
\textsc{J. J. Guadalupe, M. P\'erez and J. L. Varona,} Weak behaviour 
of Fourier-Jacobi series, \textit{to appear in Journal of Approximation 
Theory.}

\bibitem{7.}
\textsc{R. A. Hunt, B. Muckenhoupt and R. L. Wheeden}, Weighted 
norm inequalities for the conjugate function and the Hilbert 
transform, \textit{Trans. Amer. Math. Soc.} \textbf{176} (1973), 227--251.

\bibitem{8.}
\textsc{A. M\'at\'e, P. Nevai and V.Totik}, Necessary conditions 
for weighted mean convergence of Fourier series in orthogonal 
polynomials, \textit{J. Approx. Theory} \textbf{46} (1986), 314--322.

\bibitem{9.}
\textsc{B. Muckenhoupt}, Mean convergence of Jacobi series, 
\textit{Proc. Amer. Math. Soc.} \textbf{23} (1969), 306--310.

\bibitem{10.}
\textsc{B. Muckenhoupt}, Mean convergence of Hermite and Laguerre 
series~I, \textit{Trans. Amer. Math. Soc.} \textbf{147} (1970), 419--431.

\bibitem{11.}
\textsc{B. Muckenhoupt}, Mean convergence of Hermite and Laguerre 
series~II, \textit{Trans. Amer. Math. Soc.} \textbf{147} (1970), 433--460.

\bibitem{12.}
\textsc{C. J. Neugebauer}, Inserting $A_{p}$-weights, \textit{Proc. 
Amer. Math. Soc.} \textbf{87} (1983), 644--648.

\bibitem{13.}
\textsc{J. Newman and W. Rudin}, Mean convergence of orthogonal 
series, \textit{Proc. Amer. Math. Soc.} \textbf{3} (1952), 219--222.

\bibitem{14.}
\textsc{H. Pollard}, The mean convergence of orthogonal series~I,
\textit{Trans. Amer. Math. Soc.} \textbf{62} (1947), 387--403.

\bibitem{15.}
\textsc{H. Pollard}, The mean convergence of orthogonal series~II,
\textit{Trans. Amer. Math. Soc.} \textbf{63} (1948), 355--367.

\bibitem{16.}
\textsc{H. Pollard}, The mean convergence of orthogonal series~III,
\textit{Duke Math. J.} \textbf{16} (1949), 189--191.

\bibitem{17.}
\textsc{E. A. Rahmanov}, On the asymptotics of the ratio of 
the orthogonal polynomials.~II, \textit{Math. USSR Sbornik} \textbf{46} 
(1983), 105--117.

\bibitem{18.}
\textsc{G. Szeg\H{o}}, ``Orthogonal Polynomials'', 
Amer. Math. Soc. Colloq. Publ., vol. 23, Amer. Math. Soc., Providence, 
RI, 1959.

\bibitem{19.}
\textsc{G. M. Wing}, The mean convergence of orthogonal 
series, \textit{Amer. J. Math.} \textbf{72} (1950), 792--808.

\end{thebibliography}

\end{document}
